\section{Approach A: cluster, then full scan}
\label{sec:scan}
This is the reference candidate-search method. Once clusters and eligible rows are chosen,
it scores every eligible row. There is no approximation \emph{inside} that selected pool.
The work is regular and can read packed codes sequentially when filtering does not scatter them.

\fig{scan}{Approach A. The selected pool is smaller than the corpus only because of routing or filters.}{1}

\subsection{Data layout and preprocessing}
The common index from Equation~\eqref{eq:base} is sufficient. Cluster $j$ occupies the row
range between offsets $j$ and $j+1$. A row contains its packed code, and the corresponding ID
array maps the row back to the original document.

\begin{lstlisting}[caption={Candidate search for Approach A.}]
ScanSelectedClusters(q, selected_clusters, eligible, table, R):
    best = empty top-R heap
    for cluster j in selected_clusters:
        for row in eligible rows of j:
            score = ScoreDocument(code[row], table)
            best.offer(score, document_id[row])
    return best in descending score, breaking ties by ID
\end{lstlisting}

The loop is over $\Nelig$ rows, not necessarily $NP/C$. If all clusters are searched and no
filter removes documents, $\Nelig=N$. Empty filtered clusters require no document scoring,
although finding that they are empty can still have a cost.

\subsection{Step-by-step cost}
\begin{center}\small
\begin{tabularx}{\linewidth}{@{}P{.25\linewidth}YP{.23\linewidth}@{}}\toprule
Step & Work & Live temporary data\\\midrule
Locate cluster ranges & $P$ offset/range accesses & selected cluster IDs\\
Visit eligible rows & $\Nelig$ row iterations & current row/iterator\\
Score rows & $\Nelig G$ lookup terms & shared score table, accumulator\\
Keep best $R$ & $H_s(\Nelig,R)$ selection operations & top-$R$ heap\\
Return candidates & sort at most $R$ results if needed & result array or heap storage\\\bottomrule
\end{tabularx}\normalsize\end{center}

Writing $c_{\rm range}$ for range setup and $T_{\rm rank}(R)$ for final candidate ordering,
\begin{equation}
T_{A,\rm candidates}=P c_{\rm range}+\Nelig G c_{\rm lut}
                    +H_s(\Nelig,R)c_{\rm heap}+T_{\rm rank}(R).
\label{eq:scan-time}
\end{equation}
The general end-to-end latency is
\begin{equation}
\boxed{T_A=\Tcommon+T_{A,\rm candidates}+\Tfinish(R').}
\label{eq:scan-total}
\end{equation}
For the stated heap implementation, the candidate-stage time complexity is
\[
O\!\left(P+\Nelig\ceil{d/g}+\Nelig\log(R+1)+R\log(R+1)\right).
\]
The lookup term count is not a claim that each term is exactly one CPU instruction.

\subsection{What is read, and what must remain in RAM?}
The packed-code payload read for scoring is
\begin{equation}
V_{A,\rm codes}=\Nelig\,s_{\rm code}.
\end{equation}
ID reads, filter data, table accesses and result-heap accesses are additional traffic. A cache
may satisfy many of them. Thus this is a named payload count, not a complete DRAM measurement.
Sequential and scattered rows may have different effective $c_{\rm lut}$ or $c_{\rm row}$.

Index storage and candidate-stage scratch are
\begin{align}
M_{A,\rm index}&=\Mbase,\\
Q_{A,\rm search}&=\Qbase+Q_{\rm row\ iterator}+Q_{\rm scan\ workspace}.
\label{eq:scan-memory}
\end{align}
A streaming scan does not need an $N$-element score array. A matrix-based batch implementation
may choose to allocate such an array; then that allocation must be added. The same algorithmic
ranking can therefore have different peak RAM under different implementations.

\subsection{A routing miss in the four-bit example}
Partition the eight documents from Section 1 into
\[
\mathcal C_0=\{1,3,4,6\},\qquad \mathcal C_1=\{2,5,7,8\}.
\]
This is an illustrative fixed partition, not a claim about the output of a particular
clustering run. Suppose the router opens only $\mathcal C_1$. A scans scores
$1.2,0.6,0.8,-0.2$ and returns documents 2 and 7. The global top two are 1 and 2, so recall is
$1/2$. The scan did nothing wrong inside the opened cluster: document 1 was never considered.

Opening both clusters returns documents 1 and 2. This is why a scan can be exact in its selected
pool and still have imperfect global recall.

\subsection{SIMD changes a coefficient, not the score}
SIMD means one instruction operates on several data lanes. It can accelerate code extraction,
lookup accumulation, and selection. It does not by itself change the number of stored bits.
There are two useful ways to model its effect.

\textbf{Calibrated model.} Measure the implemented kernel's effective time per work unit,
$c_{\rm lut}$, for the relevant dimensions and access pattern. If that rate already includes
selection, use a combined $T_{\rm score+select}$ and do not add the heap cost a second time.

\textbf{Resource model.} If a kernel performs $U_{\rm op}$ compute units and causes $V_{\rm DRAM}$ bytes of DRAM
traffic, compute throughput $\Gamma$ and memory bandwidth $\beta$ give the lower bound
\begin{equation}
T_{\rm kernel}\geq\max\{U_{\rm op}/\Gamma,\ V_{\rm DRAM}/\beta\}.
\label{eq:resource}
\end{equation}
An effective SIMD factor $s_{\rm simd}$ can change $\Gamma$ to $s_{\rm simd}\Gamma_0$.
It cannot make the bandwidth term disappear. Random access, synchronization and fixed overhead
can put actual time above the bound. Do not add this resource estimate on top of a fitted
coefficient that already includes the same work.

For multiple queries, a float score matrix can be written as $QB^\top$, where rows of $B$ are
expanded binary codes. That exposes matrix operations, but expansion, batching delay, device
transfers and score-buffer storage have costs. A fast GPU batch is not automatically a fast
single-query request.

\begin{resultbox}{What Approach A guarantees}
For a fixed eligible pool, A returns its exact binary top $R$ under the stated tie rule.
With every eligible cluster searched, it has 100\% binary recall. With partial routing,
its global binary recall is limited by which true neighbors reach the selected pool.
No property of the scan restores documents missed by routing, quantization or an earlier filter.
\end{resultbox}
